\section{Validated primitive operations}
Real intervals are closed dyadic intervals $[l2^{-224},h2^{-224}]$ with
integer endpoints; complex intervals are Cartesian products. Addition is
exact at this scale. For multiplication take the minimum/maximum of all four
integer endpoint products, divide by $2^{224}$ and round down/up. For a
reciprocal whose interval omits zero, reverse the endpoint order and round
$2^{448}/h$ down and $2^{448}/l$ up. Negative denominators and scalar division
use the same endpoint rules. Each operation encloses all represented real
values by monotonicity; induction proves inclusion for every expression
built from them. Rectangular complex multiplication uses these real rules.
The sum of coordinate halfwidths bounds the complex distance from its exact
rational center. Binary floating arguments are never silently converted.

For $0<x<1$, the alternating arctangent sum and the endpoint obtained by
adding the next signed term enclose $\arctan x$. Sixty terms at $1/5$ and
$1/239$, together with Machin's identity
$\pi=16\arctan(1/5)-4\arctan(1/239)$, give the interval for $\pi$.
The negative coefficient reverses the second interval. The identity follows
from tangent addition and the range of the two angles; no decimal value of
$\pi$ is a trusted input.

For rational $x>0$, range-reduce $x=2^kv$, $1\le v<2$. With
$z=(v-1)/(v+1)\le1/3$,
\[
 \log v=2\sum_{j=0}^{J-1}\frac{z^{2j+1}}{2j+1}+\mathcal R,
 \qquad 0\le\mathcal R\le
 \frac{2z^{2J+1}}{(2J+1)(1-z^2)}.
\]
The geometric tail bounds each remaining denominator below by $2J+1$.
Use $J=92$, the analogous expression for $\log2$, and sign-aware interval
multiplication by $k$. This covers negative $k$ as well.

For complex $z$, repeatedly halve until its interval $\ell^1$ magnitude
bound is $A\le1/2$. The degree $60$ Taylor polynomial for $e^z$ then has remainder
at most
\[
 \frac{A^{61}}{61!}\frac1{1-A/62}.
\]
Successive absolute terms have ratio at most $A/62$, so geometric comparison
proves this bound. Add its symmetric interval to each coordinate, then square
outward the required number of times. The identity $e^{2z}=(e^z)^2$ and
interval inclusion prove the final enclosure. Caches are type-sensitive:
otherwise a cached exact argument can accidentally bypass rejection of an
equal-valued float. No implementation change is allowed during the recorded run;
the producer verifies source bytes again before emitting its certificate.

\section{Recomputing the compact growth constants}
For the larger growth disk, $c=1/2$, $r=12/5$, take $k=2$ when $a=1$
and $k=1$ when $a=0$. Both real Mellin exponents on that disk are at most
$k$. Define
\[
 Q_k(x)=\sum_{j=0}^k\frac{k!}{(k-j)!x^{j+1}},\qquad
 U_{16}(x)=(1+x/16)^{-16},\qquad \delta=3/q.
\]
Integration by parts gives
$\int_1^\infty t^ke^{-xt}\,dt=e^{-x}Q_k(x)$.
The same Gaussian split and geometric-tail comparison used below, now
with cutoff $n_0=8$, imply the exact rational majorant
\[
 \widehat M_q=2\left[
 \sum_{n=1}^{7}n^aU_{16}(\delta n^2)Q_k(\delta n^2)
 +C_aQ_k(64\delta)\frac{U_{16}(32\delta)}{1-U_{16}(8\delta)}\right],
 \quad C_0=1,\quad C_1=(1+\delta^{-1})/2.
\]
Replacing the actual $d=\pi/q$ by the smaller positive $\delta$ increases
the integral majorant. Exact fraction comparison gives
$\widehat M_3<4$, $\widehat M_4<10$, $\widehat M_5<6$.
The producer recomputes these fractions, asserts the strict comparisons and
exports them; the constants in the readout proof are thus checked rather
than supplied by a floating growth estimate.

\section{The completed-function integral and its tails}
Let $d=\pi/q$ and $\theta_q(t)=\sum_{n\ge1}\chi_q(n)n^ae^{-dn^2t}$.
The standard primitive theta transformation, followed by splitting the Mellin
integral at1, gives
\[
 \Lambda_q(s)=\int_1^\infty\theta_q(t)t^{(s+a)/2-1}\,dt+
              \int_1^\infty\theta_q(t)t^{(1-s+a)/2-1}\,dt.
\]
The second root-number coefficient is1 for these three characters: their
finite Gauss sums are $i\sqrt3$, $2i$ and $\sqrt5$, respectively.
For the last, $2\cos(2\pi/5)-2\cos(4\pi/5)=\sqrt5$; the first two are direct
conjugate-root differences. Division by $i^a\sqrt q$ gives the claimed
normalization. Exponential decay makes the two integrals entire. This classical
theta identity, not an assumption of Artin holomorphy or desired positivity,
is the analytic input of the implementation.

The actual circle nodes are enclosed by computing $e^{2\pi i j/64}$ with
the validated primitives. All input boxes fit
$D=[-1,2]+i[-4/3,4/3]$. Thus both exponents
$p=(s+a)/2-1$ and $p'=(1-s+a)/2-1$ satisfy
$|\Re p|\le2$, $|\Im p|\le2/3$. The domain guard precedes all term discards.
Retain $1\le n\le11$ and integrate $1\le t\le128$.

Put $\delta=3/q<d$, $U(x)=(1+x/128)^{-128}$ and
$Q(x)=x^{-1}+2x^{-2}+2x^{-3}$. Since $(1+x/128)^{128}\le e^x$,
$e^{-x}\le U(x)<1$ for $x>0$. Integration by parts gives
$\int_1^\infty t^2e^{-xt}\,dt=e^{-x}Q(x)$.
With $C_0=1$, $C_1=(1+\delta^{-1})/2$ the omitted-$n$ tail of both integrals is
\[
 2C_aQ(\delta12^2)
 \frac{U(\delta12^2/2)}{1-U(12\delta)}.
\]
Indeed $n^ae^{-\delta n^2/2}\le C_a$, the function $Q$ decreases, and
$(12+j)^2\ge12^2+24j$. These inequalities yield a geometric bound after
splitting the exponent in halves. The retained-$n$ time tail is at most
\[
 \sum_{n=1}^{11}2n^a128^3U(128\delta n^2)Q(128\delta n^2).
\]
It follows by substituting $t=128u$ and majorizing $t^{\Re p}$ by $t^2$.
The two tails are disjoint: omitted $n$ use the full integral, whereas
retained $n$ contribute only $t>128$.

\section{High-order integration with a certified remainder}
Split each dyadic interval $[2^j,2^{j+1}]$, $j=0,\ldots,6$, into four equal
panels. For its midpoint $m$ and halfwidth $h$, $h/m\le1/9$.
For $f(z)=e^{-\lambda z}z^p$, $\lambda=\pi n^2/q$, the principal logarithm
is analytic on $|z-m|\le\rho=m/2$. On that circle,
\[
 |f(z)|\le M_\rho=3\max\{(3m/2)^2,(m/2)^{-2},1\}
 U(\delta n^2m/2).
\]
Here $\Re z\ge m/2$, $|\arg z|\le1$ and
$e^{|\Im p|}\le e^{2/3}\le3$. The maximum handles both signs of $\Re p$.
Cauchy's coefficient estimate and cancellation of odd powers on $[-h,h]$
therefore bound the error after degree $2K$ by
\[
 \frac{2hM_\rho(h/\rho)^{2K+2}}
 {(2K+3)(1-(h/\rho)^2)}.
\]
Indeed integrate each remaining even term, replace its denominator by the
first omitted denominator, then sum a geometric series. Here $K=12$.
Multiply by $2n^a$ for both Mellin terms and sum all $28$ panels and retained $n$.
Keeping this majorant even for later discarded terms only overcounts error.

To evaluate the polynomial stably, set $u=h/m$, $q_{-1}=0$, $q_0=1$ and use
\[
 q_{j+1}=\frac{u(p-\lambda m-j)q_j-(\lambda m)u^2q_{j-1}}{j+1}.
\]
Multiplying $f'=(-\lambda+p/t)f$ by $t=m+x$ and matching coefficients proves
this recurrence, with $q_j=f^{(j)}(m)h^j/(j!f(m))$. The integrated polynomial is
$2he^{-\lambda m}m^p\sum_{j=0}^{K}q_{2j}/(2j+1)$; all its arithmetic is outward.

Fixed absolute exponential-rounding errors can be amplified by high-order
coefficients in late panels. The code therefore first bounds a complete
$(n,\text{panel})$ pair, including both powers, by
\[
 4hn^a(m+h)^2U(\delta n^2(m-h)).
\]
For real $t\ge1$ its phase has modulus1 and $t^{\Re p}\le(m+h)^2$, proving
the bound. If it is below $2^{-90}$, skip that pair entirely and add its bound
to the final radius. This is controlled truncation, not numerical underflow.
The final sample radius is the interval-sum halfwidth plus the omitted-$n$,
time, Cauchy and discard bounds. All are recorded separately and checked as
nonnegative. No high-precision special-function routine is trusted here.

\section{Reproduction and remaining interfaces}
Run the preserved source on the existing Linux Python environment with no new
packages: \texttt{python3} on \path{examples/analytic_oracle/theta_oracle.py}, with argument \texttt{--output NEW}.
It refuses to overwrite an existing experiment directory. The final-v5
experiment contains $192$ samples, three rational degree $28$ polynomials, copied
source and SHA-256 hashes. Independent validation checks rational endpoint
arithmetic, transcendental enclosures, $1160$ exact coefficient identities, high-
precision completed-function comparisons and the exported component sums.

The fixed-parameter oracle solves a concrete evaluation/certification gap.
It does not identify a positive arithmetic operator, force real-rooted
polynomials, prove the signed original-family short-row bound or supply
complete Lean4 verification. Surface/order and neural architectures may use
these certificates as independent targets, but a small learned loss is not a
substitute for the remaining arithmetic and categorical theorems.
